\documentclass[11pt]{article}

\usepackage[margin=1in]{geometry}
\usepackage{amsmath,amssymb,amsthm}
\usepackage{enumitem}
\usepackage{hyperref}
\usepackage{xcolor}
\usepackage{booktabs}

\hypersetup{colorlinks=true, linkcolor=blue, urlcolor=blue, citecolor=blue}

\title{TOAE209/TOAH209 -- Design and Analysis of Algorithms\\[4pt]
\large Week 6: Theoretical Questions}
\author{Foreign Trade University}
\date{}

\begin{document}
\maketitle

\section*{Instructions}

Part A contains 17 questions requiring written explanations, short proofs, or
recurrence-solving. Part B is a 10-question quiz (true/false and multiple choice). Attempt
every question on your own before consulting Part C (Solutions) at the end of this document.

\section*{Part A: Theory Questions}

\begin{enumerate}[label=\textbf{A\arabic*.}, leftmargin=2.2em]

\item Describe the \emph{divide-and-conquer} paradigm in general terms. Name and explain
its three steps, and write the general form of the running-time recurrence that results.

\item State the \textsc{Merge} subroutine for two sorted lists and prove that it produces a
sorted list containing exactly the elements of both inputs. State a loop invariant and use
it in your proof.

\item Prove that \textsc{Mergesort} correctly sorts any input array, by strong induction on
the array length.

\item Write down the recurrence for mergesort's running time and solve it by
\emph{unrolling} (repeated substitution). Show the intermediate steps and identify the
number of levels.

\item Solve the same mergesort recurrence $T(n) = 2T(n/2) + cn$ using a \emph{recursion
tree}. For each level $k$, give the number of nodes, the size of each node, and the total
work at that level; then sum over all levels.

\item State the \textbf{Master Theorem} in full, including all three cases and the
condition (comparison of $d$ with $\log_b a$) that selects each case.

\item For each of the following recurrences, apply the Master Theorem and give the
$\Theta$-class: (a) $T(n) = 2T(n/2) + \Theta(n)$; (b) $T(n) = T(n/2) + \Theta(1)$; (c)
$T(n) = 3T(n/2) + \Theta(n)$; (d) $T(n) = 4T(n/2) + \Theta(n)$; (e) $T(n) = 2T(n/2) +
\Theta(n^2)$. State which case applies and why.

\item For each recurrence, give the $\Theta$-class via the Master Theorem and identify the
case: (a) $T(n) = 9T(n/3) + \Theta(n)$; (b) $T(n) = T(n/2) + \Theta(n)$; (c) $T(n) =
8T(n/2) + \Theta(n^3)$; (d) $T(n) = 7T(n/2) + \Theta(n^2)$.

\item Define an \emph{inversion} of an array. How many inversions does a sorted array of
length $n$ have? How many does a strictly reverse-sorted array of length $n$ have? Justify
the reverse-sorted count.

\item Explain how counting inversions can be done in $O(n \log n)$ by modifying mergesort.
In particular, when merging sorted halves $L$ and $R$ and the algorithm emits an element of
$R$ before exhausting $L$, exactly how many inversions does that single step account for,
and why?

\item Write down and justify the recurrence for the inversion-counting algorithm. Why is it
the same order as mergesort itself?

\item State the recurrence for binary search and solve it. Then explain why binary search
is "divide and conquer with $a = 1$" and what that means for the shape of its recursion
tree (compared to mergesort's).

\item Describe the divide-and-conquer algorithm for the maximum-subarray problem. In
particular, describe the three cases the maximum subarray can fall into, and how the
\emph{cross-boundary} case is computed in $O(n)$. State and solve the recurrence.

\item Explain exponentiation by squaring. Write the recursive identity for $x^n$ (even and
odd cases), state the recurrence for the number of multiplications, solve it, and explain
why computing $x^{n/2}$ \emph{once} (rather than twice) is essential to achieving
$\Theta(\log n)$.

\item Describe \emph{quickselect} for finding the $k$-th smallest element. Explain why, with
a random pivot, the \emph{expected} running time is $\Theta(n)$, and why the Master Theorem
does not directly apply to its recurrence. What is the worst-case running time, and when
does it occur?

\item Explain the divide-and-conquer algorithm for finding a \emph{majority element} (a
value occurring more than $n/2$ times). Why does it suffice to consider only the majority
candidates of the two halves in the combine step? State and solve the recurrence.

\item Karatsuba multiplication satisfies $T(n) = 3T(n/2) + \Theta(n)$, while the naive
divide-and-conquer multiplication satisfies $T(n) = 4T(n/2) + \Theta(n)$. Solve both
recurrences with the Master Theorem and explain precisely why reducing the number of
recursive multiplications from $4$ to $3$ changes the asymptotic running time (and why
shaving the $\Theta(n)$ combine term would \emph{not} help).

\end{enumerate}

\newpage
\section*{Part B: Quiz}

\begin{enumerate}[label=\textbf{B\arabic*.}, leftmargin=2.2em]

\item \textbf{(Multiple choice)} The recurrence for mergesort is:
\begin{enumerate}[label=(\alph*)]
    \item $T(n) = T(n-1) + \Theta(n)$
    \item $T(n) = 2T(n/2) + \Theta(n)$
    \item $T(n) = 2T(n/2) + \Theta(1)$
    \item $T(n) = T(n/2) + \Theta(n)$
\end{enumerate}

\item \textbf{(True/False)} Merging two sorted lists of total length $n$ can be done in
$O(n)$ time.

\item \textbf{(Multiple choice)} By the Master Theorem, $T(n) = 2T(n/2) + \Theta(n)$ solves
to:
\begin{enumerate}[label=(\alph*)]
    \item $\Theta(n)$
    \item $\Theta(n \log n)$
    \item $\Theta(n^2)$
    \item $\Theta(\log n)$
\end{enumerate}

\item \textbf{(Multiple choice)} For $T(n) = a T(n/b) + \Theta(n^d)$, Case 1 of the Master
Theorem (where $T(n) = \Theta(n^{\log_b a})$) applies when:
\begin{enumerate}[label=(\alph*)]
    \item $d > \log_b a$
    \item $d = \log_b a$
    \item $d < \log_b a$
    \item $a = b$
\end{enumerate}

\item \textbf{(Short answer)} How many inversions does the array $[3, 2, 1]$ have? List
them.

\item \textbf{(True/False)} Counting the inversions of an $n$-element array requires
$\Theta(n^2)$ time in the worst case; no faster algorithm is possible.

\item \textbf{(Multiple choice)} The number of multiplications used by exponentiation by
squaring to compute $x^n$ is:
\begin{enumerate}[label=(\alph*)]
    \item $\Theta(n)$
    \item $\Theta(\log n)$
    \item $\Theta(n \log n)$
    \item $\Theta(\sqrt{n})$
\end{enumerate}

\item \textbf{(Multiple choice)} The \emph{expected} running time of quickselect with a
random pivot is:
\begin{enumerate}[label=(\alph*)]
    \item $\Theta(\log n)$
    \item $\Theta(n)$
    \item $\Theta(n \log n)$
    \item $\Theta(n^2)$
\end{enumerate}

\item \textbf{(True/False)} The divide-and-conquer maximum-subarray algorithm runs in
$\Theta(n \log n)$, while Kadane's algorithm solves the same problem in $\Theta(n)$.

\item \textbf{(Short answer)} Karatsuba's recurrence is $T(n) = 3T(n/2) + \Theta(n)$. Give
its $\Theta$-class (a numeric exponent to three decimal places is fine).

\end{enumerate}

\newpage
\section*{Part C: Solutions}

\subsection*{Part A}

\begin{enumerate}[label=\textbf{A\arabic*.}, leftmargin=2.2em]

\item \textbf{Divide and conquer} solves a problem by reducing it to smaller instances of
the same problem. Its three steps are: \textbf{Divide} -- split the size-$n$ input into $a$
subproblems each of size about $n/b$; \textbf{Conquer} -- solve each subproblem recursively
(stopping at small base cases); \textbf{Combine} -- merge the subproblem solutions into a
solution for the original instance, doing some additional work. The running time satisfies a
recurrence
\[
    T(n) = a\,T(n/b) + (\text{divide + combine cost}),
\]
typically written $T(n) = a\,T(n/b) + \Theta(n^d)$.

\item \textsc{Merge}$(L,R)$ keeps pointers $i,j$ into the two sorted lists and repeatedly
appends $\min(L[i], R[j])$ to the output $M$, advancing the corresponding pointer; when one
list is exhausted, it appends the remaining tail of the other. \textit{Loop invariant:} at
the start of each iteration, (i) $M$ is sorted, and (ii) every element already in $M$ is
$\le$ every element remaining in $L[i\mathbin{:}]$ and $R[j\mathbin{:}]$. \textit{Proof:}
initially $M$ is empty, so the invariant holds vacuously. Each iteration appends the smaller
of $L[i], R[j]$, which (because $L,R$ are sorted) is the smallest of all remaining elements,
hence $\le$ everything still remaining and (by the invariant) $\ge$ everything already in
$M$; so $M$ stays sorted and the invariant is preserved. When the loop ends, the remaining
tail of one list is internally sorted and all its elements are $\ge$ everything in $M$, so
appending it leaves $M$ sorted. Each input element is appended exactly once, so $M$ contains
exactly the multiset union.

\item \textit{Strong induction on $n = |A|$.} \textbf{Base case:} $n \le 1$; a list of
length $0$ or $1$ is already sorted and returned unchanged. \textbf{Inductive step:} for $n
\ge 2$, both halves have length $< n$, so by the inductive hypothesis the recursive calls
return them sorted. By A2, \textsc{Merge} of two sorted halves returns a sorted permutation
of their union, which is $A$. Hence \textsc{Mergesort}$(A)$ is correct.

\item Unrolling $T(n) = 2T(n/2) + cn$:
\[
\begin{aligned}
T(n) &= 2T(n/2) + cn = 2\bigl(2T(n/4) + c\tfrac n2\bigr) + cn = 4T(n/4) + 2cn \\
     &= 8T(n/8) + 3cn = \cdots = 2^k T(n/2^k) + k\,cn.
\end{aligned}
\]
The recursion stops when $n/2^k = 1$, i.e.\ $k = \log_2 n$, giving $2^k T(1) = n\cdot O(1) =
O(n)$ plus accumulated combine cost $k\,cn = cn\log_2 n$. Therefore $T(n) = \Theta(n \log
n)$. There are $\log_2 n + 1$ levels.

\item Recursion tree for $T(n) = 2T(n/2) + cn$: level $k$ (root is level $0$) has $2^k$
nodes, each of size $n/2^k$, each doing $c(n/2^k)$ combine work. Total work at level $k$ is
$2^k \cdot c(n/2^k) = cn$ -- the \emph{same} at every level. There are $\log_2 n + 1$
levels, so the total is $cn(\log_2 n + 1) = \Theta(n \log n)$.

\item \textbf{Master Theorem.} For $T(n) = a\,T(n/b) + \Theta(n^d)$ with $a \ge 1$, $b > 1$,
$d \ge 0$:
\begin{itemize}
    \item \textbf{Case 1} ($d < \log_b a$, i.e.\ $a > b^d$): $T(n) = \Theta(n^{\log_b a})$ --
    leaves dominate.
    \item \textbf{Case 2} ($d = \log_b a$, i.e.\ $a = b^d$): $T(n) = \Theta(n^d \log n)$ --
    every level does equal work.
    \item \textbf{Case 3} ($d > \log_b a$, i.e.\ $a < b^d$): $T(n) = \Theta(n^d)$ -- the root
    dominates.
\end{itemize}

\item (a) $a=2,b=2,d=1$, $\log_2 2 = 1 = d$ $\Rightarrow$ Case 2, $\Theta(n \log n)$. (b)
$a=1,b=2,d=0$, $\log_2 1 = 0 = d$ $\Rightarrow$ Case 2, $\Theta(\log n)$. (c) $a=3,b=2,d=1$,
$\log_2 3 \approx 1.585 > 1$ $\Rightarrow$ Case 1, $\Theta(n^{\log_2 3}) \approx
\Theta(n^{1.585})$. (d) $a=4,b=2,d=1$, $\log_2 4 = 2 > 1$ $\Rightarrow$ Case 1,
$\Theta(n^2)$. (e) $a=2,b=2,d=2$, $\log_2 2 = 1 < 2$ $\Rightarrow$ Case 3, $\Theta(n^2)$.

\item (a) $a=9,b=3,d=1$, $\log_3 9 = 2 > 1$ $\Rightarrow$ Case 1, $\Theta(n^2)$. (b)
$a=1,b=2,d=1$, $\log_2 1 = 0 < 1$ $\Rightarrow$ Case 3, $\Theta(n)$. (c) $a=8,b=2,d=3$,
$\log_2 8 = 3 = d$ $\Rightarrow$ Case 2, $\Theta(n^3 \log n)$. (d) $a=7,b=2,d=2$, $\log_2 7
\approx 2.807 > 2$ $\Rightarrow$ Case 1, $\Theta(n^{\log_2 7}) \approx \Theta(n^{2.807})$.

\item An \textbf{inversion} of an array $A$ is a pair of indices $(i,j)$ with $i < j$ but
$A[i] > A[j]$. A sorted (non-decreasing) array has \textbf{0} inversions (no pair is out of
order). A strictly reverse-sorted array of length $n$ has \textbf{$\binom{n}{2} =
\frac{n(n-1)}{2}$} inversions: \emph{every} one of the $\binom{n}{2}$ pairs $i < j$ has
$A[i] > A[j]$, which is the maximum possible.

\item Run mergesort; whenever \textsc{Merge} emits an element $R[j]$ from the right half
while $L$ still has $|L| - i$ elements remaining (i.e.\ $R[j] < L[i]$, so the algorithm
takes from the right), those $|L| - i$ remaining left elements are each \emph{to the left}
of $R[j]$ in the original array yet \emph{larger} than $R[j]$ -- so each forms exactly one
inversion with $R[j]$. Thus a single such step accounts for exactly $|L| - i$ inversions.
Summed across all merges, this counts every cross-inversion exactly once; combined with the
inversions counted recursively inside each half, the total is the array's inversion count.

\item $T(n) = 2T(n/2) + \Theta(n)$, identical to mergesort's, because the only addition to
\textsc{Merge} is incrementing a counter by $|L| - i$ at each right-emission step -- $O(1)$
extra work per merge step, leaving the per-merge cost $\Theta(n)$ unchanged. By the Master
Theorem (Case 2), $T(n) = \Theta(n \log n)$.

\item Binary search: $T(n) = T(n/2) + \Theta(1)$. By the Master Theorem with $a=1$, $b=2$,
$d=0$ ($\log_2 1 = 0 = d$, Case 2), $T(n) = \Theta(\log n)$. It is "divide and conquer with
$a=1$" because, although the input is split into two halves, the algorithm \emph{recurses
into only one} of them -- so the recursion tree is a single \emph{path} of length $\log_2 n$
(not a branching tree as in mergesort), and the total work is just $\Theta(1)$ per level
times $\log_2 n$ levels.

\item Split $A$ at the midpoint $m$. The maximum-sum subarray is one of: (i) entirely in the
left half, (ii) entirely in the right half, or (iii) \emph{crossing} the midpoint. Cases (i)
and (ii) are solved by recursion. For (iii), the best crossing subarray is the best suffix
of the left half (computed by scanning leftward from $m$ keeping a running maximum) joined to
the best prefix of the right half (scanning rightward) -- $O(n)$ total. Return the max of the
three. The recurrence is $T(n) = 2T(n/2) + \Theta(n) = \Theta(n \log n)$ (Master Theorem,
Case 2).

\item Exponentiation by squaring uses
\[
    x^n = \begin{cases} 1 & n=0,\\ (x^{n/2})^2 & n \text{ even},\\ x\cdot(x^{(n-1)/2})^2 & n \text{ odd}.\end{cases}
\]
Letting $M(n)$ be the number of multiplications, $M(n) = M(\lfloor n/2 \rfloor) + O(1)$
(one or two multiplications to square and possibly multiply by $x$). By the Master Theorem
($a=1,b=2,d=0$, Case 2), $M(n) = \Theta(\log n)$. The key is that $x^{n/2}$ is computed
\emph{once} and squared; if instead we wrote $x^n = x^{n/2}\cdot x^{n/2}$ and computed each
factor separately, the recurrence would become $M(n) = 2M(n/2) + O(1) = \Theta(n)$ -- no
savings over naive repeated multiplication.

\item \textbf{Quickselect}: choose a pivot $p$, partition $A$ into elements $<p$, $=p$, and
$>p$; if the $k$-th smallest lies in the $<p$ part, recurse there; if it is among the $=p$
part, return $p$; otherwise recurse on the $>p$ part with $k$ adjusted by the number of
discarded smaller elements. With a \emph{random} pivot, with constant probability the pivot
discards a constant fraction of the array, so the expected work satisfies $T(n) \le T(cn) +
O(n)$ for some constant $c < 1$ (e.g.\ $c = 3/4$), a decreasing geometric series summing to
$O(n)$. The Master Theorem does not apply because (a) only \emph{one} subproblem is solved
and (b) its size is a constant \emph{fraction} of $n$ rather than a fixed $n/b$ for both
branches. The worst case is $\Theta(n^2)$, occurring when the pivot is repeatedly the
largest or smallest remaining element (e.g.\ a sorted input with a last-element pivot), so
each call shrinks the problem by only one.

\item \textbf{Majority element via D\&C}: split into halves; recursively find each half's
majority candidate. A value that is a majority of the \emph{whole} (occurs $> n/2$ times)
must be a majority of at least one half -- because if it were a strict minority in
\emph{both} halves (occurring $\le$ half the time in each), it would occur $\le n/2$ times
overall, a contradiction. So the only candidates are the two halves' majority candidates;
the combine step counts each candidate's occurrences across the whole range (an $O(n)$
scan) and returns whichever exceeds $n/2$. Recurrence: $T(n) = 2T(n/2) + \Theta(n) =
\Theta(n \log n)$.

\item Karatsuba: $a=3,b=2,d=1$; $\log_2 3 \approx 1.585 > 1$ $\Rightarrow$ Case 1, $T(n) =
\Theta(n^{\log_2 3}) \approx \Theta(n^{1.585})$. Naive D\&C: $a=4,b=2,d=1$; $\log_2 4 = 2 >
1$ $\Rightarrow$ Case 1, $T(n) = \Theta(n^2)$. Both are in Case 1, so the running time is
$\Theta(n^{\log_b a})$, governed \emph{entirely} by $a$ (the number of recursive
multiplications): dropping $a$ from $4$ to $3$ drops the exponent from $2$ to $\log_2 3
\approx 1.585$. Shaving the $\Theta(n)$ combine term would not help, because in Case 1 the
combine term is asymptotically dominated by the leaf work $n^{\log_b a}$ -- the running time
does not depend on $d$ at all once $d < \log_b a$.

\end{enumerate}

\subsection*{Part B}

\begin{enumerate}[label=\textbf{B\arabic*.}, leftmargin=2.2em]

\item \textbf{(b)} $T(n) = 2T(n/2) + \Theta(n)$: two recursive sorts of half-size halves,
plus a linear merge.

\item \textbf{True.} Each comparison emits one element, so the total number of
comparisons/appends is at most the total length $n$; merging is $\Theta(n)$.

\item \textbf{(b)} $\Theta(n \log n)$. ($a=b=2$, $d=1$, $\log_2 2 = 1 = d$, Master Theorem
Case 2.)

\item \textbf{(c)} $d < \log_b a$ (equivalently $a > b^d$): leaves dominate, $T(n) =
\Theta(n^{\log_b a})$.

\item \textbf{3 inversions:} $(0,1)$ since $3>2$, $(0,2)$ since $3>1$, and $(1,2)$ since
$2>1$. This is the maximum $\binom{3}{2} = 3$ for a length-3 array.

\item \textbf{False.} Counting inversions can be done in $\Theta(n \log n)$ via the modified
mergesort; the $\Theta(n^2)$ brute force is not optimal.

\item \textbf{(b)} $\Theta(\log n)$ -- each step halves the exponent.

\item \textbf{(b)} $\Theta(n)$ expected (with a random pivot). The worst case is
$\Theta(n^2)$.

\item \textbf{True.} The D\&C version is $\Theta(n \log n)$ (recurrence $2T(n/2) +
\Theta(n)$); Kadane's single-pass dynamic program is $\Theta(n)$.

\item $\Theta(n^{\log_2 3}) \approx \Theta(n^{1.585})$.

\end{enumerate}

\end{document}
